\section{Generalization Analysis}

\begin{tcolorbox}
\textbf{Q1:} From the perspective of stability, how does the introduction of the weight decay fundamentally reshape the stability bound of the learning algorithm? To maintain optimal system performance, what dependencies must other hyperparameters exhibit? 
\end{tcolorbox}

In fact, \citet{10.5555/2986916.2987033} have offered a preliminary answer assuming a linear network, showing that weight decay significantly enhances generalization by suppressing irrelevant weight components and mitigating the impact of noise. However, this conclusion cannot be extended to parameter-dependent theoretical studies (e.g., involving $\lambda, \eta, T$), precluding further analysis of optimal parameter settings. Consequently, we investigate the theoretical properties of the weight decay algorithm based on stability.

%Establishing a generalization error bound for certain algorithm and using it to find variable dependencies is a primary means of studying generalization ability. 

\subsection{Stability and Generalization Definition}
This section derives generalization bounds for various algorithms via \emph{uniform stability}. We start by outlining the theoretical connection between generalization and stability.
First, we introduce the following notions of \emph{uniform stability} and \emph{the expected generalization error}.
\begin{definition}[$\epsilon$-Uniformly Stable]
A randomized algorithm $A$ is $\epsilon$-uniformly stable if for all data sets $S, S^{\prime} \in Z^{n}$ such that $S$ and $S^{\prime}$ differ in at most one example, we have
    \begin{equation}\label{E-stab}
     \mathop{sup}\limits_{z\in Z}\left\{\mathbb{E}_{A}\left[F(A_S;z)-F(A_{S^{\prime}};z)\right] \right\} \leq \epsilon.
    \end{equation}
\end{definition}

We define the expected generalization error of a model $\theta = A_S$, obtaining via a randomized algorithm $A$, as
\begin{equation}\label{D-gen}
\mathbb{E}_{S,A}\left[R_{S}\left[A_S\right]-R_\mathcal{D}\left[A_S\right]\right]. 
\end{equation}

By leveraging the concept of \emph{uniform stability}, \citet{hardt2016train} provide the generalization bound for the SGD algorithm based on the following theorem.
\begin{theorem}{\rm (Generalization in Expectation, \citet[Theorem 2.2]{hardt2016train})}\label{E-gen}
Let $A$ be $\epsilon$-uniformly stable. Then,
    \begin{equation}
     \vert\mathbb{E}_{S,A}\left[R_{S}\left[A_S\right]-R_\mathcal{D}\left[A_S\right]\right]\vert \leq \epsilon.
    \end{equation}
\end{theorem}

%In SGD, the main sources of randomness include the model itself and the random selection of data. Therefore, we need to take expectations on both $A$ and $S$. 
\begin{remark}
Here, the expectation is taken only over the internal randomness of $A$. Theorem \ref{E-gen} implies that an algorithm with uniform stability exhibits small generalization error. Therefore, the generalization bound can be determined through controlling the algorithm's uniform stability.   
\end{remark}

%We recall the important theorem that uniform stability implies \emph{generalization in expectation} \cite{hardt2016train}. The proof is based on \citet[Lemma 7]{bousquet2002stability} and very similar to \citet[Lemma 11]{shalev2010learnability}.

Stability investigates the impact of input perturbations on model outputs. To characterize this, we analyze the update states of the original and perturbed sequences. Let $S$ and $S^{\prime}$ denote two datasets differing by a single sample. After running the optimizer for $T$ steps on $S$ and $S^{\prime}$, the resulting parameters are denoted as $\theta_T$ and $\theta_T^{\prime}$. Therefore, $\mathbb{E}\vert F(\theta_T;z)-F(\theta_T^{\prime};z)\vert$ reflects the algorithm's stability under perturbations. We employ the Lipschitz property of the loss function and then obtain
\begin{equation}\label{convex-basic}
\mathbb{E}\vert F(\theta_T;z)-F(\theta_T^{\prime};z)\vert \leq G \mathbb{E}\Vert \theta_T -\theta^{\prime}_T \Vert
 \end{equation}
for all $\theta_T$ and $\theta_T^{\prime}$.
This implies that we will further transform the problem of obtaining the generalization bound via the stability bound into one of controlling parameter differences. 

Next, we consider the algorithm's expansion properties and establish stability bounds based on them.

\begin{lemma}{\rm (\citet[Lemma 3.6]{hardt2016train})}
\label{lemma3.6}
Assume that the function $F$ is $L$-smooth and convex. Then for any $\eta \leq \frac{2}{L}$, we have $\Vert \theta_{T+1}-\theta_{T+1}^{\prime} \Vert \leq \Vert \theta_{T}-\theta_{T}^{\prime}\Vert.$
\end{lemma}

\begin{remark}
Lemma \ref{lemma3.6} indicates that when a convex function satisfies smoothness, and the step size is sufficiently small, perturbations exert minimal influence on the optimizer, and gradient updates exhibit non-expansive behavior. The proof is detailed in Appendix \ref{pro-lemma}. Further results can be found in papers \citep{hardt2016train,xiao2022stability}.    
\end{remark}

%The next section will derive generalization bounds for different algorithms based on the properties of gradient updates.

\subsection{Generalization bounds for algorithms}

\begin{lemma}{\rm (\citet[Theorem 3.8]{hardt2016train})}\label{thm:sgd}
Assume that the loss function $F(\theta;z)$ is convex, Lipschitz and $L$-smooth for all given $z\in\mathcal{Z}$ with sizes $n$. Suppose we run SGD with learning rate $\eta \leq \frac{2}{L}$ for $T$ steps, where each step samples $z_i$ from $\mathcal{Z}$ uniformly with replacement. Then, SGD has uniform stability of
\begin{equation}
  \epsilon_{gen} = \mathbb{E}\vert F(\theta_T;z)-F(\theta^{\prime}_T;z)\vert \leq \frac{2\eta G^2 T}{n}.
 \end{equation}
\end{lemma}

\begin{remark}
Lemma \ref{thm:sgd} establishes a generalization upper bound dependent on the learning rate $\eta$ and the number of iterations $T$, facilitating the analysis of optimal parameter settings. Furthermore, given that \citet{pmlr-v180-zhang22b} identified this upper bound $\mathcal{O}(T/n)$ as also being the lower bound, it provides a theoretical foundation for comparing the generalization improvements of decay algorithms.    
\end{remark}

\begin{theorem}\label{thm:sgd-WD}
Assume that the loss function $F(w;z)$ is convex, Lipschitz and $L$-smooth for all given $z\in\mathcal{Z}$ with sizes $n$. Suppose we run SGD-DW with weight decay coefficient $\lambda \in (0,1)$ and learning rate $\eta \leq \frac{2}{2\lambda +L}$ for $T$ steps, where each step samples $z_i$ from $\mathcal{Z}$ uniformly with replacement. Then, SGD-DW has uniform stability of 
\begin{equation}
  \epsilon_{gen} = \mathbb{E}\vert F(\theta_T;z)-F(\theta^{\prime}_T;z)\vert \leq \frac{2G^2}{\lambda n}.
 \end{equation}
\end{theorem}

\begin{remark}
Theorem \ref{thm:sgd-WD} presents a bound $\mathcal{O}(1/n)$, which is tighter than the bound $\mathcal{O}(T/n)$ in Lemma \ref{thm:sgd}. Unlike the linear growth with respect to iterations $T$ in SGD, SGD-WD achieves a smaller constant upper bound, indicating that weight decay significantly improves generalization. This indicates that weight decay is not merely a regularization technique; rather, it alters the optimization dynamics, steering the algorithm toward regions with improved stability. This insight underlies our connection to SWA, which similarly seeks flatter, more stable regions of the landscape. 
\end{remark}

\begin{remark}\label{remark 2}
The tighter bound of SGD-WD is primarily attributed to the weight decay coefficient $\lambda$, which induces the $(1-\lambda\eta)$-expansive property in the algorithm, denoted $\Vert \theta_t-\theta^{\prime}_t\Vert \leq (1-\lambda\eta)\Vert \theta_{t-1} - \theta^{\prime}_{t-1}\Vert$. This aligns with the theoretical properties of SGD under the strong convexity assumption \cite{hardt2016train}, yielding an equivalent upper bound. However, a key distinction is that SGD-WD necessitates a smaller learning rate to satisfy the expansive condition under convexity, which is essential for convergence. For the detailed proof, please refer to appendix \ref{sec:app-thm4.1}.    
\end{remark}

\paragraph{In summary,} 

\begin{theorem}\label{thm:sgdm}
Assume that the loss function $F(w;z)$ is $\mu$-strongly convex, Lipschitz and $L$-smooth for all given $z\in\mathcal{Z}$ with sizes $n$. Suppose we run SGDM with learning rate $\eta \leq \frac{2}{L(1-\beta)}$ for $T$ steps, where each step samples $z_i$ from $\mathcal{Z}$ uniformly with replacement. Then, SGDM has uniform stability of 
\begin{equation}
  \epsilon_{gen} = \mathbb{E}\vert F(\theta_T;z)-F(\theta^{\prime}_T;z)\vert \leq \frac{2\tau_{\beta}\eta G^2 T}{n},
 \end{equation}
where $\tau_{\beta}\in(0,1)$ is a constant, and as $T\rightarrow\infty$, $\tau_{\beta}\rightarrow 1$. 
\end{theorem}

\begin{remark}
The Theorem \ref{thm:sgdm} shows that, under $\mu$-strongly convex assumptions, SGDM admits a generalization bound $\mathcal{O}(\tau_{\beta}\eta T/n)$, with the improvement reflected in a reduced constant $\tau_\beta$ dependent on $\beta$, highlighting the stabilizing effect of momentum.
In fact, $\tau_\beta=[1-(1-\beta)^t]\beta+(1-\beta)$, and as $T\rightarrow\infty$, $\tau_{\beta}\rightarrow 1$. We retain it to account for momentum's influence.    
\end{remark}

% \begin{remark}
    
% \end{remark}

\begin{theorem}\label{thm:sgdm-WD}
Assume that the loss function $F(w;z)$ is $\mu$-strongly convex, Lipschitz and $L$-smooth for all given $z\in\mathcal{Z}$ with sizes $n$. Suppose we run SGDM-WD with weight decay coefficient $\lambda \in (0,1)$ and learning rate $\eta \leq \frac{2}{2\lambda + L\Gamma_{\beta}}$ for $T$ steps, where $\Gamma_{\beta}\in(0,1)$ is a constant dependent on the momentum coefficient $\beta$ and each step samples $z_i$ from $\mathcal{Z}$ uniformly with replacement. Then, SGDM-WD has uniform stability of
\begin{equation}
  \epsilon_{gen} = \mathbb{E}\vert F(\theta_T;z)-F(\theta^{\prime}_T;z)\vert \leq \frac{2\tau_{\beta} G^2 }{\lambda n},
 \end{equation}
 where $\tau_{\beta}\in(0,1)$ is a constant, and as $T\rightarrow\infty$, $\tau_{\beta}\rightarrow 1$. 
\end{theorem}

\begin{remark}
Compared to Theorem \ref{thm:sgdm}, Theorem \ref{thm:sgdm-WD} shows that SGDM-WD achieves a tight bound $\mathcal{O}(1/n)$ that is independent of iteration $T$. Furthermore, this bound improves the result of Theorem \ref{thm:sgd-WD}, particularly in terms of the factor $\tau_{\beta}$. For the introduced momentum $\beta$, this fully demonstrates the dominant role of weight decay $\lambda$.
\end{remark}

\begin{remark}\label{remark 1}
While weight decay tightens the generalization bound, it simultaneously imposes a stricter constraint on the learning rate, specifically $\eta \leq \frac{2}{2\lambda + L\Gamma_{\beta}}$ to ensure the non-expansive property of SGD-WD under convex assumptions. The term $2\lambda$ in the denominator reveals an intrinsic trade-off: enhancing regularization strength (larger $\lambda$) necessitates a proportional reduction in the step size $\eta$ to prevent the contractive mechanism from destabilizing the trajectory. For clarity, complex momentum-related constants are subsumed into $\Gamma_{\beta} \in (0,1)$, as they affect the precise stability boundary but do not alter the structural coupling between $\lambda$ and $\eta$. The proof for explicit factors $\tau_{\beta}$ and $\Gamma_{\beta}$ are detailed in Appendix \ref{sec:app-thm:sgdm-wd}. 
\end{remark}

While weight decay tightens the generalization bound, it simultaneously imposes a stricter constraint on the learning rate, specifically $\eta \leq \frac{2}{2\lambda + L\Gamma_{\beta}}$, to maintain the non-expansive property. The term $2\lambda$ in the denominator reveals an intrinsic trade-off: enhancing regularization strength (larger $\lambda$) necessitates a proportional reduction in the step size $\eta$ to prevent the contractive mechanism from destabilizing the trajectory. For clarity, complex momentum-related constants are subsumed into $\Gamma_{\beta} \in (0,1)$, as they affect the precise stability boundary but do not alter the structural coupling between $\lambda$ and $\eta$.

\begin{figure}[t]
    \centering
    \includegraphics[width=0.85\linewidth]{figures/performance.jpg}
    \vspace{-0.4cm}
    \caption{The comparison of SGD, SGD-WD, SGDM, and SGDM-WD in optimal accuracy and corresponding learning rates.}
    \label{fig:performance}
    \vspace{-0.4cm}
\end{figure}

\textbf{Answer to Q1:} As demonstrated by Lemma \ref{thm:sgd} and Theorems \ref{thm:sgd-WD}, \ref{thm:sgdm} and \ref{thm:sgdm-WD}, weight decay not only enhances the generalization capability of base optimizers but also imposes new requirements on learning rate setting. Specifically, for the basic optimizer, the introduction of $\lambda$ significantly reduces the conditions of learning rate under optimal performance, which is also empirically verified, as shown in Figure \ref{fig:performance}. This prompts us to rethink the coupling between $\lambda$ and $\eta$.


